% ---
% id: EX-CH01-005
% titre: Nombres parfaits
% chapitre_id: 1
% chapitre_nom: Nombres entiers & divisibilité
% annee_scolaire: 9H
% difficulte: 2
% tags: [divisibilite, diviseurs]
% auteur: Bussy D.
% ---

\begin{exercice}{EX-CH01-005}{Nombres parfaits}

\begin{mdframed}[backgroundcolor=vlmgray, innerleftmargin=10pt, innerrightmargin=10pt, innertopmargin=6pt, innerbottommargin=6pt, skipabove=-5pt, skipbelow=5pt]
Un \textbf{nombre parfait} est un nombre naturel non nul qui est égal à la somme de ses diviseurs stricts (ses diviseurs sauf lui-même).
\end{mdframed}


\bigskip
Parmi les nombres suivants, lesquels sont parfaits ?

\bigskip
\sol{\begin{tabularx}{\linewidth}{*{5}{>{\centering\arraybackslash}X}}
6 & 28 & 496 & 875 & 8128\\
\end{tabularx}}%
{
On détermine les diviseurs stricts de chaque nombre, puis on vérifie si leur somme est égale au nombre.

\vspace{20pt}
\def\arraystretch{1.2}
\begin{tabularx}{\linewidth}{lXlc}
$D^*(6)$    & $= \{1, 2, 3\}$ & $\sum = 6$    & $\checkmark$ \\
$D^*(28)$   & $= \{1, 2, 4, 7, 14\}$                                             & $\sum = 28$  & $\checkmark$ \\
$D^*(496)$  & $= \{1, 2, 4, 8, 16, 31, 62, 124, 248\}$                           & $\sum = 496$  & $\checkmark$ \\
$D^*(875)$  & $= \{1, 5, 7, 25, 35, 125, 175\}$                                  & $\sum = 373 \neq 875$ & $\times$ \\
$D^*(8128)$ & $= \{1, 2, 4, 8, 16, 32, 64, 127, 254, 508, 1016, 2032, 4064\}$   & $\sum = 8128$ & $\checkmark$ \\
\end{tabularx}
\def\arraystretch{1}

\vspace{20pt}
Les nombres parfaits parmi ceux proposés sont donc : 6, 28, 496 et 8128.
}
\end{exercice}